\section{Pretraining}
\paragraph{Purpose}
Train the model from scratch on a very large corpus (web pages, books, code) with the
next-token objective of the language modeling chapter. It is self-supervised, and the result
is a \tB{base model}: it continues a text very well, but it is not yet an assistant that
follows instructions.

\paragraph{Ingredients}
\begin{itemize}
    \item \textbf{data}: collection, filtering and \uB{deduplication}. The quality of the
        data matters as much as its quantity;
    \item the texts are tokenized and cut into windows of $n_{ctx}$ tokens, batches contain
        millions of tokens;
    \item \textbf{AdamW} with warmup, cosine decay and gradient clipping (see the
        optimizers section), in reduced precision (bfloat16);
    \item the model and the batches are spread over thousands of accelerators
        (data, tensor and pipeline parallelism).
\end{itemize}

\paragraph{Compute}
For a model with $N$ parameters trained on $D$ tokens, the forward pass costs about $2ND$
FLOPs and the backward pass about twice more:
\begin{center}
    \frB{$C\approx 6\,N\,D$ FLOPs}
\end{center}
Example: GPT-3, $N=175\cdot10^{9}$ and $D\approx300\cdot10^{9}$, gives
$C\approx3\cdot10^{23}$ FLOPs.

\section{Scaling Laws}
\paragraph{Observation} (Kaplan et al., 2020) The test loss decreases as a \uB{power law}
of the number of parameters $N$, of the data $D$ and of the compute $C$, smoothly over many
orders of magnitude, as long as none of them is the bottleneck. It makes the loss of a large
training run \tB{predictable from small runs}.

\paragraph{Compute-optimal training} (Hoffmann et al., ``Chinchilla'', 2022) For a fixed
budget $C=6ND$, there is a best split between $N$ and $D$. They found that the parameters and
the tokens should \tB{grow together}, with about \uB{20 tokens per parameter}. Previous large
models were undertrained (GPT-3: about 2 tokens per parameter): a 70B model trained on 1.4T
tokens outperformed larger models trained with the same compute.

\paragraph{Caveats}
\begin{itemize}
    \item the optimum above only accounts for the \emph{training} cost. A model that will
        answer a very large number of requests costs more at inference when it is bigger, so
        it is worth training a smaller model on more tokens than the Chinchilla rule;
    \item the laws predict the loss, not directly the performance on a given downstream task;
    \item the data is finite, repeating it many times brings less and less.
\end{itemize}
